\chapter{Conclusiones y recomendaciones}
\label{cap:conclusiones}

\section{Conclusiones}

\subsection{Sobre la pregunta de investigación}

La pregunta era qué indicadores cuantitativos de cobertura y organización de la roca pueden
derivarse de forma reproducible de las máscaras de AI4Mars, y cuáles son sus límites de validez
frente a diferencias en la anotación y al juicio humano. La respuesta tiene tres partes: la
cobertura sí puede derivarse de forma reproducible; el conteo de rocas individuales, no, y ese
resultado negativo es en sí mismo un hallazgo; y el propio conjunto de datos presenta
inconsistencias que delimitan lo que cualquiera de los dos indicadores puede afirmar.

La \textbf{cobertura de roca visible} puede derivarse de forma reproducible para prácticamente
todo el subconjunto. Pudo calcularse en \cNCobCalc{} de las \cNEscenas{} escenas
(\cPctCobCalc); \cNCobPos{} escenas (\cPctCobPos{} del total) presentaron cobertura de roca
mayor que cero, y entre ellas la mediana sobre píxeles etiquetados fue del \cCobMedVal. No se
encontró evidencia de que sus valores altos se expliquen por una menor fracción de escena
etiquetada. Su límite de validez es la fuente de la anotación: con máscaras de especialista el
valor sería algo menor, de modo que las coberturas describen el terreno tal como lo registra la
anotación colaborativa.

El \textbf{conteo de rocas} es el resultado negativo del trabajo: \textbf{contar rocas
individuales a partir de estas máscaras no es sencillo}, y el procedimiento no lo consigue más
que en un sentido acotado: produce un conteo reproducible de las instancias
geométricamente separables dentro de la clase roca grande, pero ese conteo no puede
interpretarse como estimación válida del número total de rocas visibles en la escena. La
evaluación exploratoria con un observador sugiere un subconteo sistemático dentro de la propia
región anotada, debido sobre todo a regiones que no contienen la información geométrica
necesaria para separar instancias; y como la clase no cubre toda la roca de la escena, ni
siquiera un conteo exacto de esa región equivaldría al número de rocas visibles.

Las \textbf{inconsistencias del conjunto de datos} completan la respuesta. Las escenas con
máscara de especialista son mucho menos rocosas que las de entrenamiento; la diferencia de
criterio entre voluntarios y especialistas, si existe, es pequeña e incierta; la roca grande es
la clase que menos consenso reúne
incluso entre especialistas, y la anotación de la roca grande es parcial y de
granularidad variable. Estos rasgos, que la Tabla~\ref{tab:inconsistencias} reúne, explican buena
parte de los límites de los dos indicadores.

\subsection{Sobre los objetivos específicos}

\textbf{E1} produjo la cobertura en dos versiones complementarias —sobre píxeles etiquetados y
sobre la imagen completa—, acompañadas de la fracción etiquetada que permite juzgar la
diferencia. Las conclusiones principales se mantienen con cualquiera de las dos versiones.

\textbf{E2} produjo el conteo para las \cFlagOk{} escenas de bandera \texttt{ok}, con un total de
\cNRocas{} rocas y una distribución tamaño--frecuencia decreciente, coherente en forma con la
literatura de abundancia de rocas. En \cBandaCero{} de esas escenas (\cBandaCeroPct\,\%) los
filtros de área y forma descartan toda la roca grande y el conteo es cero. El criterio de
prominencia redujo la fragmentación de la primera versión, pero la reconstrucción de la
calibración muestra que lo hizo rebajando el conteo de forma general y no solo en las escenas
sobresegmentadas, y no la eliminó: la división de aguas subdivide alguna región en el
\cPctSubdiv{} de las escenas que cuentan rocas. La población de E2 excluye las \cNNullBig{}
escenas casi sin etiquetar que contienen roca grande, por ser las que concentran los artefactos
de anotación.

\textbf{E3} caracterizó la composición del terreno, definió una tipología de escenas y mostró la
alternancia entre tramos rocosos y tramos dominados por el suelo a lo largo de la secuencia de
adquisición. Esa variación es temporal y no espacial: el orden de adquisición no equivale a la
distancia recorrida. La variación entre tipos de escena se describe con la propia tipología,
que se define a partir de la composición y, por tanto, de la cobertura. Tampoco fue posible
comparar los dos ojos del par estéreo, porque el subconjunto es casi enteramente monocular, con
\cNOjoIzq{} imágenes del ojo izquierdo frente a \cNOjoDer{} del derecho.

\textbf{E4} entregó el conjunto de recursos reproducibles: código modular con parámetros
explícitos, guiones de ejecución que regeneran todas las cifras, tablas y figuras de este
documento, registro de versiones y un archivo tabular con una fila por imagen y veintiocho
columnas de identificación, indicadores y elegibilidad.

\textbf{E5} estableció los límites de validez mediante tres contrastes —fuente de la anotación,
juicio de un observador y segmentador entrenado, a los que se suma el control del denominador
de E1— y atribuyó los errores del conteo a sus causas: en el \cPctSinBig{} de las escenas no hay
roca grande que contar, en el \cPctFiltradas{} los filtros la descartan, y entre las escenas que
cuentan rocas la subdivisión afecta al \cPctSubdiv. De ahí proceden los hallazgos que se resumen
a continuación.

\subsection{Sobre las hipótesis}

La Tabla~\ref{tab:contraste-hipotesis} recoge el dictamen de cada hipótesis. \textbf{H1} se
sostiene: la cobertura es reproducible y no está gobernada por la fracción etiquetada, aunque
persiste una dependencia no lineal débil. \textbf{H2} se rechaza, en una evaluación
exploratoria: el conteo subestima de forma sistemática frente al observador. \textbf{H3} se
rechaza con reservas: hay un efecto de la fuente de la anotación, pequeño e incierto, muy
inferior a la diferencia bruta entre fuentes, que procede sobre todo de las imágenes.
\textbf{H4} se sostiene: cobertura y conteo aportan información distinta, sin ser
independientes.

\subsection{Sobre la relación entre los indicadores}

Los indicadores presentan una asociación lineal prácticamente nula en las escenas analizadas
($r = \cDepCCPearson$) y tampoco una asociación monótona. Sí muestran, en cambio, una
dependencia débil y no lineal: la correlación de distancia y la información mutua la detectan
($p \cDepCCDcorP$ y $p \cDepCCIMP$), con forma de U invertida, en parte inducida por la
propia definición de los indicadores. La magnitud es pequeña: los tramos de cobertura explican
el \cEtaConteo{} de la varianza del conteo. La conclusión que se sostiene es que los dos
indicadores aportan información distinta sobre la anotación —cuánta superficie se etiquetó como
roca y en cuántas regiones separables se etiquetó la roca grande— y que reportarlos por separado
está justificado, aunque la relación entre ambos no sea nula.

\subsection{Sobre la fuente de la anotación}

La comparación directa entre máscaras colaborativas y de experto muestra una diferencia grande:
la cobertura mediana de las escenas con roca pasa del \cCobMedVal{} al \cGoldCobMedPos. Esa
diferencia, sin embargo, \textbf{no puede atribuirse sin más a la anotación}, porque las dos
fuentes cubren imágenes distintas. Aplicando a una muestra de cada fuente un mismo segmentador
como instrumento común, el \cDescImgPct{} de la diferencia de cobertura media entre las muestras
(\cDescTotal{} puntos) se asocia con la composición de las imágenes y \cDescAnot{} puntos
porcentuales (IC~95\,\% \cDescAnotIC) con la fuente de las etiquetas.

De ahí se siguen dos conclusiones. La primera es que \textbf{las máscaras de experto de AI4Mars
para esta cámara cubren escenas mucho menos rocosas} que las de entrenamiento, una propiedad del
conjunto que afecta a cualquier comparación entre ambos. La segunda es que la componente
atribuible a la anotación es pequeña e incierta: su estimación es positiva en todas las
variantes —entre \cDescAnot{} y \cFVDescAnot{} puntos según la muestra y el instrumento, en el
sentido de que la anotación colaborativa asigna algo más de roca que la de especialista—, pero
cuando se tiene en cuenta que las imágenes de experto se concentran en \cTempNSolesExp{} soles,
su intervalo en la muestra principal incluye el cero \cTempAnotIC. Los
resultados son compatibles con la hipótesis de un sesgo de saliencia en la anotación
colaborativa, pero el diseño del estudio no permite identificar causalmente ese mecanismo. El
mecanismo alternativo de un denominador reducido por suelo sin etiquetar no encuentra apoyo,
porque la fracción etiquetada no difiere apreciablemente entre las dos fuentes
(\cFracValidMed{} frente a \cFracValidGold, $p = \cFracValidP$).

\subsection{Sobre el segmentador entrenado}

Un segmentador DeepLabV3 entrenado con las máscaras colaborativas, con un reparto por bloques
temporales que deja la prueba a más de un sol del entrenamiento, reproduce la cobertura de esas
máscaras con un IoU medio de \cSegMIoU{} en la muestra equilibrada de prueba y de \cSegNatMIoU{}
en su distribución natural. Un reparto anterior al azar por imagen daba cifras prácticamente
iguales, de modo que la proximidad entre imágenes no las inflaba. Frente a las máscaras de
experto el IoU medio baja a \cExpMIoU{} y el modelo sobreestima la cobertura en \cBAExpMedia{}
puntos de media \cBAExpMediaIC, lo que es compatible con que haya heredado el criterio de la
anotación con que se entrenó, aunque estos datos no permiten demostrarlo.

Frente al experto, el modelo presenta por tanto un sesgo positivo. El error absoluto medio es de
\cBAExpMAE{} puntos, los límites de acuerdo abarcan unos sesenta y cinco puntos porcentuales y el
error depende del tramo de cobertura, con sobreestimación en las coberturas intermedias y
subestimación en las de cobertura total. El modelo describe conjuntos de escenas, no escenas
individuales, y estima cobertura, no conteo.

\subsection{Sobre el conteo y el juicio humano}

La evaluación exploratoria por bandas sobre \cVN{} escenas, con un observador, arrojó un acuerdo
exacto del \cVAcuerdo, con kappa de Cohen de \cVKappa{} y kappa ponderado de \cVKappaPond: un
acuerdo escaso. El desacuerdo es asimétrico —el procedimiento queda por debajo del observador en
\cVAbajo{} escenas y por encima en \cVArriba— y no depende del ajuste de los parámetros: en las
\cSenN{} combinaciones ensayadas el procedimiento subcuenta en al menos \cSenAbajoMin{} escenas.

La evidencia apunta a que la causa principal del subconteo está en la máscara más que en el
algoritmo. Cuando un polígono envuelve un campo de bloques contiguos, la máscara no contiene la
información geométrica que permitiría separarlos, y la división de aguas no tiene dónde cortar.
El procedimiento separa de forma reproducible lo que la geometría de la máscara permite
separar, con una subdivisión residual en parte de las escenas
(Sección~\ref{sec:diagnostico}); la máscara permite separar mucho menos de lo que un observador
distingue. Aparte de ese desacuerdo, la anotación de roca grande no marca todas las rocas de una
escena sino algunas —en \cVNoExhN{} de las \cVN{} escenas representa menos del
\SI{20}{\percent} de la roca etiquetada—, de modo que el conteo tampoco representaría la escena
completa aunque coincidiera con el observador dentro de la región.

\subsection{Sobre el origen de los límites del conteo}

La evidencia indica que la principal limitación observada para el conteo de instancias proviene
de la representación semántica y de la granularidad de las etiquetas, más que de una falta de
resolución evidente de las imágenes. El lecho rocoso aporta el \cParteBed{} de la roca
etiquetada y la roca grande solo el \cParteBig; y sobre las mismas \cNGold{} imágenes, la
proporción de escenas con roca grande cae del \cGoldUnoBig{} al \cGoldTresBig{} al endurecer el
acuerdo exigido a los especialistas: solo se conserva el \cConsBig{} de sus píxeles, frente al
\cConsBed{} del lecho rocoso, lo que indica que la roca grande es la clase más difícil de
delimitar también para ellos. No es posible, con todo, aislar la taxonomía como causa única: intervienen
también las reglas de consenso, la perspectiva, la escala, la oclusión y la parametrización.

Los ensayos que usan la imagen para separar bloques dentro de la región anotada no mejoran el
acuerdo con el observador de forma demostrable una vez corregido el número de métodos probados.
El enfoque híbrido queda como prueba de concepto.

\subsection{Sobre las inconsistencias del conjunto de datos}

El diagnóstico del conjunto de datos es el resultado más transferible del trabajo. Seis
inconsistencias —cinco documentadas con evidencia cuantitativa y una ilustrada con ejemplos en
la Sección~\ref{sec:inconsistencias}—
condicionan su uso como medida del terreno y como referencia de evaluación: la composición
distinta de las escenas de entrenamiento y de especialista, la posible diferencia de criterio
entre voluntarios y especialistas, el escaso consenso sobre la roca grande, la anotación no exhaustiva de
la roca grande, la granularidad variable del trazado y las escenas casi sin etiquetar. Ninguna
invalida el conjunto para el propósito con que se creó, pero todas deben tenerse en cuenta antes
de leer sus máscaras como medidas.

\section{Aportes}

\paragraph{Un diagnóstico de los límites de validez de las máscaras de AI4Mars.} Es el aporte
principal. El trabajo establece qué puede y qué no puede inferirse de esas máscaras: la
cobertura se deriva de forma reproducible y no está gobernada por la fracción etiquetada; el
conteo de instancias queda limitado sobre todo por la representación semántica y la granularidad
de las etiquetas, junto con otros factores; y la comparación entre
fuentes de anotación está dominada por una diferencia de composición entre las muestras de
imágenes que, hasta donde llega la revisión realizada, no se había documentado. Esto último es
pertinente para todo trabajo que emplee las máscaras de experto como referencia de evaluación.

\paragraph{Un procedimiento y un conjunto de resultados.} Un flujo documentado, auditable y de
bajo costo computacional que deriva, para \cNEscenas{} escenas, una fila de indicadores por
imagen con veintiocho columnas, ejecutable en procesador convencional sin conjunto de entrenamiento. Todas las cifras,
tablas y figuras se regeneran con guiones a partir de los resultados, y el conjunto es
reutilizable con independencia del procedimiento que lo generó.

\paragraph{Un diseño de evaluación transferible.} Dos elementos del diseño son aplicables a
otros conjuntos anotados: medir la relación entre indicadores con estadísticos capaces de
detectar dependencia no lineal, y no solo con Pearson; y separar el efecto de la anotación del
efecto de las imágenes mediante un instrumento común cuando las fuentes no comparten imágenes.

\paragraph{Un resultado negativo documentado.} Mostrar con datos por qué un conteo
geométricamente razonable no reproduce el juicio humano sobre estas máscaras, y que el problema
no se resuelve recalibrando umbrales, delimita qué cabe esperar de AI4Mars para tareas de conteo
de instancias.

\paragraph{Una extensión aplicada.} Las reglas heurísticas de priorización y la aplicación de
consulta, documentadas en el Anexo~\ref{app:extension}, muestran cómo usar los indicadores con
un criterio explícito. Son un demostrador y no una contribución central, y no están calibradas
para decisiones de navegación.

\section{Recomendaciones}

\subsection{Para el uso de los resultados de este trabajo}

Las coberturas deben emplearse para \textbf{comparar escenas del mismo conjunto}, no como
estimaciones absolutas de la abundancia de roca, y teniendo presente que una anotación de
especialista produciría valores algo menores.

El conteo debe acompañarse siempre de su bandera de calidad y leerse como número de bloques
delimitados en la clase roca grande, no como número de rocas de la escena. Un cero puede
significar ausencia de rocas o ausencia de la etiqueta que permite contarlas, situaciones que el
conjunto de resultados distingue.

Los indicadores de tamaño relativo no deben leerse literalmente en escenas de fracción
etiquetada inferior a $0{,}05$ —bandera \texttt{mostly\_null}—, y conviene tratarlos con
cautela por debajo de $0{,}20$, el umbral con que se excluyeron escenas del entrenamiento del
segmentador. La variación a lo largo de la secuencia de adquisición no debe leerse
como variación a lo largo del trayecto.

\subsection{Para trabajos futuros}

\paragraph{Ampliar la validación humana.} Es la prioridad. La validación se hizo con un único
observador, el propio autor, sobre una muestra estratificada que no es representativa del
subconjunto. Convendría repetirla con al menos dos observadores independientes del autor,
idealmente tres, sobre una muestra representativa y con el mismo protocolo —imágenes
anonimizadas, orden aleatorio, resultado automático oculto—, reportando primero el acuerdo
entre personas, con un estadístico como el alfa de Krippendorff, y después el acuerdo entre el
procedimiento y cada una de ellas. Hacer dos preguntas por escena, una sobre la región anotada y
otra sobre la escena completa, permitiría además cuantificar cuánto deja fuera la anotación no
exhaustiva.

\paragraph{Evaluar el conteo basado en la imagen.} El relieve de sombras es el que más se acerca
al observador, sin superar la corrección por comparaciones múltiples. Su evaluación requiere una
muestra independiente de la usada aquí, varios observadores y un control explícito de la
sobresegmentación que la textura de la imagen introduce.

\paragraph{Probar una taxonomía con bloques explícitos.} La taxonomía geológica que AI4Mars
incluye para \textit{Perseverance} distingue roca suelta y guijarros. Extender el análisis a esa
misión permitiría poner a prueba la hipótesis de que una taxonomía con representación más
explícita de los bloques mejora la estimación de instancias. Más allá del inventario descriptivo
del Anexo~\ref{app:vetas}, no se aborda aquí porque implica otro rover, otra cámara, otra
distribución de datos y nuevas validaciones.

\paragraph{Recuperar la escala métrica.} La cámara de navegación es un par estéreo, y el archivo
planetario de la NASA distribuye para sus imágenes productos de distancia y de coordenadas
derivados del par. Combinarlos con las máscaras, aunque estas solo existan para el ojo
izquierdo, permitiría expresar los tamaños en unidades métricas y comparar las distribuciones
tamaño--frecuencia en magnitud, y no solo en forma, con la literatura de abundancia de
rocas.

\paragraph{Estimar el indicador sin anotación humana.} El segmentador aproxima la cobertura a
partir de la imagen sola. Validarlo contra varias referencias y sobre otras partes de la misión
permitiría calcular el indicador sobre imágenes sin máscara. Ese uso
requeriría, antes de cualquier aplicación a la planificación, una validación específica que este
trabajo no ofrece.

\subsection{Para quienes evalúan modelos con AI4Mars}

Las máscaras de experto de la cámara de navegación de \textit{Curiosity} cubren escenas mucho
menos rocosas que las de entrenamiento. Una caída del desempeño al pasar de unas a otras mezcla,
por tanto, el cambio de anotación con el cambio de composición, y no debe atribuirse solo a la
primera. Se recomienda reportar el desempeño frente a ambas referencias, por clase, describir la
composición de cada conjunto de evaluación y, como precaución, repartir los datos agrupando por
secuencia de adquisición para que imágenes casi idénticas no queden a ambos lados del reparto;
en este trabajo el reparto al azar no infló las cifras, pero no hay garantía de que ocurra lo
mismo con otros modelos o conjuntos.

\subsection{Para quienes mantienen el conjunto de datos}

Tres medidas reducirían las inconsistencias documentadas: describir cómo se eligieron las
imágenes de la muestra de especialista, cuya composición difiere tanto de la de entrenamiento;
publicar máscaras colaborativas de esas mismas imágenes, para que las dos fuentes puedan
compararse sobre el mismo material; y anotar la roca grande por instancias, con un polígono por
bloque, si se quiere que el conjunto sirva para contar rocas.
